% Part: many-valued-logic
% Chapter: syntax-and-semantics
% Section: connectives

\documentclass[../../../include/open-logic-section]{subfiles}

\begin{document}

\olfileid[fr]{mvl}{syn}{con}

\olsection{Langages et connecteurs}

La logique propositionnelle classique,
comme beaucoup d'autres logiques,
dispose de constantes propositionnelles
et de connecteurs. Nous prenons, par
exemple, les symboles primitifs suivants :
\begin{enumerate}
\tagitem{prvFalse}{La constante propositionnelle
  de la fausseté~$\lfalse$.}{}
\tagitem{prvTrue}{La constante propositionnelle
  de la vérité~$\ltrue$.}{}
\item Les connecteurs logiques :
  \startycommalist
  \iftag{prvNot}{\ycomma $\lnot$ (négation)}{}%
  \iftag{prvAnd}{\ycomma $\land$ (conjonction)}{}%
  \iftag{prvOr}{\ycomma $\lor$ (disjonction)}{}%
  \iftag{prvIf}{\ycomma $\lif$ (conditionnel)}{}%
  \iftag{prvIff}{\ycomma $\liff$ (biconditionnel)}{}%
\end{enumerate}
\iftag{defNot,defOr,defAnd,defIf,defIff,defTrue,defFalse,defEx,defAll}{%
Outre les connecteurs primitifs, nous
employons aussi des symboles définis
comme abréviations, tels que
\startycommalist
  \iftag{defNot}{\ycomma $\lnot$ (négation)}{}%
  \iftag{defAnd}{\ycomma $\land$ (conjonction)}{}%
  \iftag{defOr}{\ycomma $\lor$ (disjonction)}{}%
  \iftag{defIf}{\ycomma $\lif$ (conditionnel)}{}%
  \iftag{defIff}{\ycomma $\liff$ (biconditionnel)}{}%
  \iftag{defFalse}{\ycomma $\lfalse$ (fausseté)}{}%
  \iftag{defTrue}{\ycomma $\ltrue$ (vérité)}{}.}{}

Les mêmes connecteurs peuvent figurer
dans les logiques plurivalentes. Il
est toutefois parfois utile d'inclure
dans un même langage plusieurs versions,
par exemple, de la conjonction. Il
faut alors des symboles distincts pour
les distinguer. Certaines logiques
plurivalentes ont aussi des connecteurs
sans équivalent classique. Nous adopterons
donc une notion de langage plus générale.

\begin{defn}
Un \emph{langage propositionnel} est
formé d'un ensemble~$\Lang L$ de
\emph{connecteurs}. Chaque connecteur
$\star$ possède une \emph{arité};
s'il est d'arité~$n$, on dit qu'il
est \emph{$n$-aire}. Les connecteurs
d'arité~$0$ sont aussi appelés
\emph{constantes}; ceux d'arité~$1$
sont \emph{unaires}, et ceux d'arité~$2$
\emph{binaires}.
\end{defn}

\begin{ex}
Le langage propositionnel usuel
$\Lang L_0$ comprend les connecteurs
suivants, avec leur arité :
$\lfalse$~($0$),
$\lnot$~($1$),
$\land$~($2$),
$\lor$~($2$) et
$\lif$~($2$).
La plupart des logiques étudiées ici
utilisent ce langage. Certaines,
par tradition ou par convention,
emploient d'autres symboles : dans
la logique produit, par exemple,
la conjonction se note souvent
$\odot$ plutôt que~$\land$.
Il est parfois commode d'ajouter
un opérateur nouveau, tel que
l'opérateur de détermination
$\triangle$, d'arité~$1$.
\end{ex}

\end{document}
